\ficha{Tao e Salomão (2020), A concepção do padrão-ouro no século XIX}{taosalomao2020}

\begin{identificacao}
TAO, Matheus Itiro de Castro; SALOMÃO, Ivan Colangelo. Um conceito, diversas
formas: a concepção do padrão-ouro na centralidade do capitalismo no século
XIX. In: ENCONTRO DE PÓS-GRADUAÇÃO EM HISTÓRIA ECONÔMICA, 10.; CONFERÊNCIA
INTERNACIONAL DE HISTÓRIA ECONÔMICA, 8., 2020, evento virtual. \emph{Anais}.
ABPHE, 2020. 18 p.

Comunicação de evento, história do pensamento econômico, 18 páginas, em
português. Lida integralmente no PDF dos anais. Texto não revisado: a seção 6
vem depois da conclusão e para no meio de uma frase (p. 17).
\end{identificacao}

\bloco{Problema e tese}
O trabalho pergunta o que sustenta teoricamente a defesa do padrão-ouro nos
clássicos e responde que a sustenta a teoria quantitativa do dinheiro, com uma
tensão interna que o texto expõe. Se o valor do dinheiro é regulado pela
relação entre oferta e demanda, e não pelas condições de produção, o dinheiro
não precisa ser mercadoria, e o ouro não é exigência lógica do arranjo
quantitativista. A tese é que o vínculo com o ouro é conveniência prática e que
a defesa do padrão-ouro constitui um projeto de estabilidade do valor do
dinheiro e do nível de preços.

\bloco{Argumento}
\begin{enumerate}[leftmargin=*,nosep]
\item Apoiado em Germer (1999), o texto separa dois sentidos do conceito. Como
conceito histórico, padrão-ouro nomeia a configuração de um sistema monetário
em um período. Como conceito teórico, nomeia um sistema idealmente pensado em
que o ouro circula como dinheiro e a função monetária se esgota na circulação.
Os autores adotam o ângulo nacional e deixam fora do recorte o sistema
internacional de 1880 a 1914 (p. 3).

\item A asserção fundante vem de Hume, para quem os preços são proporcionais à
abundância de dinheiro (p. 4). Daí a neutralidade, definida como incapacidade
do dinheiro de afetar produto, emprego e juros (nota 9, p. 5). Da mesma
premissa segue que o dinheiro ``não tem de ser necessariamente uma mercadoria
(ouro)'' (p. 4). O problema montado é por que os quantitativistas defendem o
ouro se a teoria deles o dispensa, e a resposta é a conversibilidade, que
aciona o autoajuste da oferta de dinheiro e dos preços (p. 5).

\item Padrão exclusivo contra bimetalismo. Com dois metais em curso legal a
valor fixo, a medida-padrão passa de um ao outro conforme o valor relativo, e o
metal desmonetizado é fundido (Ricardo, p. 9). Mill acrescenta que o
bimetalismo tende à depreciação do padrão, pois os pagamentos migram para o
metal que menos subiu ou que mais caiu (p. 9). O monometalismo dá maior fixidez
ao valor do dinheiro.

\item Hume e Ricardo divergem sobre o papel-moeda. Hume vê no crédito uma
simulação de dinheiro que expulsa os metais preciosos (p. 13). Ricardo afirma
que a moeda está no estado mais perfeito quando é toda de papel de valor
equivalente ao ouro que representa (p. 14) e que o papel nem precisa ser
pagável em espécie, bastando que a quantidade seja regulada pelo valor do metal
adotado como padrão (p. 15). A lei bancária de Peel, de 1844, converte o
princípio em legislação (p. 16).
\end{enumerate}

\chave{o padrão-ouro cauciona a estabilidade do nível de preços, já que o
automatismo que lhe é inerente não anui a possibilidade do dinheiro se pôr
acima ou abaixo do nível natural de produção do país}{p. 12}

\chave{a vinculação do padrão de preços ao ouro não emana de uma suposta
necessidade lógica do sistema teórico, mas de uma questão de conveniência
prática}{p. 16}

\bloco{Conceitos e definições operacionais}
Padrão-ouro em sentido teórico: sistema monetário nacional em que o ouro é base
e meio circulante (p. 3). Medida-padrão de valor, ou padrão de preços: a
mercadoria única que mede valores, distinta do padrão monetário imposto por
lei, que pode ser bimetálico (nota 20, p. 10). Conversibilidade: obrigação de
converter notas em ouro cunhado ou em lingotes, tratada como freio à emissão e
não como fim em si (p. 5). Princípio da limitação da quantidade: o postulado
quantitativista aplicado à emissão, para Ricardo o aspecto mais importante do
papel-moeda (p. 14). Metalismo prático, na acepção schumpeteriana: defender uma
mercadoria como dinheiro sob o suposto de que as funções monetárias se limitam
a padrão monetário e meio circulante (p. 17).

\bloco{Evidência e método}
História do pensamento econômico por leitura textual de Hume, Ricardo, Mill,
Smith e Marx, todos em tradução brasileira de edições de divulgação. O apoio
secundário é Germer (1999) e Fonseca e Mollo (2012). Não há dados, recorte
empírico nem contrafactual, e nenhuma fonte primária brasileira é lida.

\bloco{Agentes e vertentes}
Do lado quantitativista, Hume como origem, Ricardo como proeminente intelectual
do bulionismo e inspirador da Currency School, além de Mill, Peel e McCulloch
(p. 2 e p. 5). Smith apoia a origem do dinheiro e o papel-moeda bem organizado.
A Banking School e os antibulionistas comparecem só em nota, remetidos a
Fonseca e Mollo (2012), com a velocidade de circulação e a não neutralidade
(nota 11, p. 6). Marx e Engels entram pela descrição da lei de 1844 e pela tese
oposta, a de que o dinheiro é mercadoria por necessidade lógica (nota 28, p.
16). No Brasil, os autores nomeiam Rodrigues Torres, Torres Homem e Jerônimo
Teixeira Júnior como herdeiros de Hume e Ricardo (p. 2 e p. 17).

\bloco{Críticas e limites}
Comunicação de evento, não artigo arbitrado, e inacabada: a seção sobre os
metalistas brasileiros tem quatro parágrafos e para no meio de uma frase (p.
17). O que ela afirma, que a transposição do receituário inglês não foi passiva
porque a economia brasileira se amparava em bases agrário-escravagistas, é
enunciado de programa, não resultado demonstrado. A leitura de Ricardo como
matriz da Currency School não é consensual, e Gontijo (2014) o coloca na
linhagem da lei do refluxo, divergência que o trabalho não discute. Ao
restringir o conceito ao plano nacional, o texto nada diz sobre regras do jogo,
credibilidade, assimetria centro-periferia, nem sobre arranjos de
conversibilidade limitada como a Caixa de Conversão.

\begin{dialogo}
\item Autoridade dos clássicos. Procurar menção nominal a Hume, Ricardo, Mill,
Peel e à lei bancária de 1844 nos artigos sobre a Caixa. Testa se em 1906 a
autoridade inglesa ainda é invocada diretamente, como faziam os metalistas
imperiais, ou se o argumento já circula sem citar a origem.

\item Ouro como conveniência e não como virtude do metal. Procurar o argumento
de que o ouro serve por dar fixidez ao valor do mil-réis, com vocabulário de
época como ``estabilidade da moeda'' e ``padrão fixo''. Testa a tese central do
trabalho no material brasileiro e ajuda a separar quem defende a estabilidade a
taxa nova de quem defende a valorização.

\item Regra de quantidade contra conversão efetiva. Ricardo admite papel não
pagável em espécie desde que a quantidade seja regulada pelo valor do ouro (p.
15). Procurar se os defensores da Caixa argumentam pelo teto de emissão e pela
proporção com o lastro, ou se exigem a conversão no guichê e a custódia do
metal. Testa qual das duas versões da doutrina chegou a 1906.

\item Padrão de preços fixado por lei. O trabalho distingue o padrão monetário
imposto pela lei da medida-padrão de valor efetiva (nota 20, p. 10). Procurar
como os jornais tratam a taxa de 15 pence diante do par legal de 27 pence: como
conveniência prática, como depreciação do padrão, ou como quebra da medida de
valor.
\end{dialogo}

\usonocapitulo{Parte III, como apoio. O texto não trata da Caixa de Conversão
nem do debate de 1906, e a única passagem sobre o Brasil está inacabada. Serve
a duas afirmações. Primeira, que a doutrina importada pelos metalistas
imperiais tem conteúdo preciso e reconstruível, ouro como padrão exclusivo,
conversibilidade como freio de emissão e limitação da quantidade, e que é esse
conteúdo que se deve procurar nos agentes de 1906. Segunda, que a defesa do
ouro é subordinada a um objetivo, a estabilidade do valor do dinheiro, o que dá
critério para distinguir estabilidade de valorização. Complementa Fonseca e
Mollo (2012), que dá o mapa das duas margens, ao descer ao interior da margem
metalista, e não substitui Gontijo (2014), cujo objeto é o mecanismo de ajuste.}
